% ============================================================
%  SSNS Project Proposal V3 — Overleaf-compatible LaTeX
%  Compact single-column academic formatting, target: ~4 pages
%  All content verbatim from project-proposal-v3.tex
% ============================================================
\documentclass[10pt,a4paper]{article}

% ---------- Packages ----------
\usepackage[margin=0.75in]{geometry}
\usepackage[utf8]{inputenc}
\usepackage[T1]{fontenc}
\usepackage{microtype}          % Better text fill / fewer overfull hboxes
\usepackage{booktabs}
\usepackage{array}
\usepackage{tabularx}
\usepackage{enumitem}
\usepackage{float}
\usepackage{cite}
\usepackage[compact]{titlesec}
\usepackage[skip=1pt,font=small]{caption}
\usepackage[hidelinks,
            colorlinks=true,
            linkcolor=black,
            citecolor=black,
            urlcolor=blue,
            pdftitle={Revised Project Proposal (V2)},
            pdfauthor={M.Sc. High Integrity Systems}
           ]{hyperref}

% Tighten lists globally — preserves all content, reduces vertical space
\setlist[itemize]{noitemsep, topsep=0pt, parsep=0pt, partopsep=0pt, leftmargin=*}
\setlist[enumerate]{noitemsep, topsep=0pt, parsep=0pt, partopsep=0pt, leftmargin=*}

% Tighter section spacing
\titlespacing{\section}{0pt}{*1.0}{*0.5}
\titlespacing{\subsection}{0pt}{*0.5}{*0.2}

% Tighter paragraph spacing
\setlength{\parskip}{0pt}
\setlength{\parindent}{1.2em}

% Tighter float spacing
\setlength{\floatsep}{4pt}
\setlength{\intextsep}{4pt}
\setlength{\textfloatsep}{4pt}

% ============================================================
\begin{document}

\title{\vspace{-2em}Thread-Based Intelligent Indoor Environmental Monitoring System\\
\large Smart Sensor Network System --- Revised Project Proposal (V2)\vspace{-0.5em}}

\author{%
Haris Tauseef Khan, Mobeen Anwar, Muhammad Faris Mufti,\\
Muhammad Inayat Hussain, Muhammad Saleem, Sohail Asghar Minhas
\vspace{0.5em}
  \\
  M.Sc.\ High Integrity Systems - Summer Semester 2026
}
\date{\vspace{-2em}}

\maketitle

% ============================================================
\begin{abstract}
\normalsize
\noindent This proposal presents a scalable IoT platform for intelligent indoor environmental
monitoring. The proposed system integrates OpenThread (IPv6) sensor nodes based on the
Nordic nRF52840 microcontroller with a containerized Gateway-Tier Fog backend. By
applying unsupervised machine learning---specifically Isolation Forest-based anomaly
detection and sliding window linear regression---across a multi-dimensional telemetry
stream, the platform moves beyond conventional single-parameter threshold alerting.
Furthermore, a lightweight Explainable AI (XAI) layer is integrated to interpret
complex, multi-variable environmental anomalies. This architecture aims to provide
operators with transparent, context-aware insights while maintaining implementation
feasibility on edge devices, ultimately delivering a reproducible and intelligent
telemetry analytics system.
\end{abstract}

% ============================================================
\section{Introduction}

Indoor air quality directly affects human health, comfort, and occupational safety.
Elevated concentrations of particulate matter, carbon dioxide, and volatile organic
compounds (VOCs) are associated with adverse respiratory and cognitive effects in
indoor environments \cite{who2021}. Despite well-established risks, many built
environments still rely on periodic manual inspection or rigid threshold-based
monitoring systems. Such approaches often lack adaptive multi-variable analysis and
may struggle to generalize across dynamic IoT and industrial monitoring environments
\cite{ortega2025,wu2022}.

This project proposes a scalable IoT platform designed to support multi-node
environmental monitoring utilizing low-power Thread sensor nodes. The system processes
telemetry using both deterministic safety rules and unsupervised machine learning. It
is designed to identify subtle multi-variable environmental anomalies and provide early
warning indicators for hazardous trends before they breach critical safety thresholds.

% ============================================================
\section{Literature Review \& Research Gap}

Recent advancements in IoT environmental monitoring have explored machine-learning
approaches for anomaly detection in resource-constrained and Industrial IoT environments.
Ortega et al.\ \cite{ortega2025} presented a lightweight Isolation Forest-based anomaly
detection approach for edge devices, while Wu et al.\ \cite{wu2022} discussed GNN-based
anomaly detection methods for Industrial Internet of Things (IIoT) systems. However,
many existing frameworks rely heavily on cloud infrastructure, introducing latency and
privacy concerns, or suffer from the ``Cold Start Problem'' when deployed in new
environments \cite{chhabra2025}. Furthermore, while dynamic sliding window techniques
have been explored for predictive accuracy in industrial systems \cite{guha2025}, their
integration into low-power indoor monitoring networks remains limited.

\textbf{Research Gap:} There is a distinct lack of localized and adaptive edge-native
IoT monitoring systems that combine unsupervised micro-climate anomaly detection with
low-power mesh networking and fog-assisted processing. This project addresses this gap
by coupling a Thread-based CoAP sensor network with a localized Gateway-Tier Fog
backend, integrating anomaly detection, trend analysis, and Explainable AI into a
cohesive, edge-deployable system \cite{chaturvedi2026}.

% ============================================================
\section{Proposed System Architecture}

\subsection{Environmental Parameters \& Sensor Mapping}

The system employs a suite of complementary sensors, measuring distinct environmental
parameters mapped to international health guidelines (ASHRAE, WHO, EPA). Sensor readings
are transmitted via CoAP over Thread.

\begin{table}[htbp]
\centering
\setlength{\tabcolsep}{4pt}
\renewcommand{\arraystretch}{1.0}
\caption{Environmental Parameters \& Sensor Mapping}
\label{tab:sensors}
\small
\begin{tabularx}{\textwidth}{@{} l l l >{\raggedright\arraybackslash}X @{}}
\toprule
\textbf{Measurement} & \textbf{Hardware} & \textbf{Typical Indoor Range} & \textbf{Health Guideline Threshold} \\
\midrule
CO\textsubscript{2}  & Sensirion SCD41   & 400 -- 1000~ppm          & $>$ 1500~ppm (ASHRAE 62.1) \\
Humidity             & Sensirion SCD41   & 30\% -- 60\%               & $>$ 70\% or $<$ 20\% (ASHRAE 55) \\
Temperature          & Microchip MCP9808 & 19\textdegree C -- 28\textdegree C  & $>$ 35\textdegree C or $<$ 10\textdegree C (ASHRAE 55) \\
PM 2.5 / PM 10       & Sensirion SPS30   & 0 -- 45~\textmu g/m\textsuperscript{3} & $>$ 35~\textmu g/m\textsuperscript{3} (WHO / EPA) \\
TVOC                 & Adafruit CCS811  & 0 -- 500~ppb             & $>$ 1000~ppb (WELL v2) \\
\bottomrule
\end{tabularx}
\end{table}

\subsection{Architecture Layers (Gateway-Tier Fog)}

The system logically decouples microservices on a local Gateway-Tier Fog appliance to
reduce latency and minimize external internet dependency.

\begin{table}[htbp]
\centering
\setlength{\tabcolsep}{4pt}
\renewcommand{\arraystretch}{1.0}
\caption{Architecture Layers (Gateway-Tier Fog)}
\label{tab:arch}
\small
\begin{tabularx}{\textwidth}{@{} >{\raggedright\arraybackslash}p{2.2cm} >{\raggedright\arraybackslash}p{2.5cm} >{\raggedright\arraybackslash}p{3.5cm} >{\raggedright\arraybackslash}X @{}}
\toprule
\textbf{Layer} & \textbf{Component} & \textbf{Technology} & \textbf{Role} \\
\midrule
Hardware Node  & nRF52840 DK        & ARM Cortex-M4, OpenThread     & Sensor acquisition \& CoAP transmission. \\
\midrule
Edge Gateway   & Serial Gateway     & Python 3.12, pyserial         & Serial-to-HTTP bridge, telemetry ingestion. \\
\midrule
Backend API    & Spring Boot        & Java 21, Spring Boot 3.3      & Telemetry routing, evaluation, and persistence. \\
\midrule
ML Engine      & ML Service         & Python, FastAPI, scikit-learn & Isolation Forest, Trend Analysis, \& XAI execution. \\
\midrule
Database       & Timeseries DB      & MongoDB 7                     & Telemetry and historical context persistence. \\
\midrule
Presentation   & Dashboard          & React, Vite                   & Visualization and alert interpretation. \\
\bottomrule
\end{tabularx}
\end{table}

\subsection{Wireless Protocol Justification (OpenThread \& CoAP)}

OpenThread (IPv6) was selected for the hardware nodes to address the strict power and
reliability constraints of indoor monitoring:

\begin{itemize}
  \item \textbf{Energy Efficiency:} Thread operates as a low-power wireless mesh network,
    allowing battery-operated sensor nodes to run continuously without excessive power draw.
  \item \textbf{Self-Healing Mesh Topology:} Thread natively forms a self-healing mesh,
    allowing packets to route automatically around physical obstacles (e.g., concrete walls)
    without requiring a direct line-of-sight to the gateway \cite{khattak2023}.
  \item \textbf{Data Reliability (CoAP):} The system utilizes CoAP over Thread, supporting
    \texttt{CONFIRMABLE} message delivery. This mitigates packet loss and enhances the
    robustness of telemetry ingestion, a strict requirement for accurate machine learning
    evaluation \cite{coap2014}.
  \item \textbf{Scalability:} For multi-floor deployments, OpenThread nodes act as
    low-power routers, bouncing telemetry to a central gateway, thereby reducing
    infrastructure overhead.
\end{itemize}

% ============================================================
\section{Intelligent Analysis \& Alerting}

\subsection{Machine Learning Pipeline}

The system avoids heavy, cloud-dependent Deep Learning models in favor of lightweight,
Edge-native algorithms suitable for Fog Computing:

\begin{itemize}
  \item \textbf{Anomaly Detection (Isolation Forest):} Analyzes multi-dimensional sensor
    data to generate a dynamic, rolling baseline for individual micro-climates. Isolation
    Forest is utilized due to its linear time complexity, minimal memory footprint, and
    unsupervised nature \cite{ortega2025,wu2022}.
  \item \textbf{Trend Analysis (Sliding Window Linear Regression):} Provides early
    warning indicators for primary environmental parameters and pollutant concentrations
    (e.g., CO\textsubscript{2}, PM2.5, TVOC, Temperature). By capturing short-term local
    trends using a sliding window, the system provides lightweight predictive support
    suitable for resource-constrained edge environments. Prior IIoT research has shown
    that many deep-learning-based sliding-window forecasting approaches introduce higher
    computational overhead and runtime complexity in low-latency deployments \cite{guha2025}.
  \item \textbf{Explainable AI (XAI) \& Interpretability:} To prevent opaque alerting,
    the ML engine implements a lightweight XAI layer. During an anomaly, the system
    calculates a baseline deviation ranking, providing interpretable contextual feedback
    directly to the operator to improve situational awareness.
\end{itemize}

\subsection{Layered Alerting Model}

The platform employs a four-tier alerting model:

\begin{enumerate}
  \item \textbf{Deterministic Safety Alerts:} Synchronous notifications triggered when
    metrics breach established health guidelines.
  \item \textbf{Behavioral Anomaly Alerts:} Asynchronous alerts triggered when
    multi-variable patterns deviate from the established micro-climate baseline.
  \item \textbf{Early Warning Indicators:} Proactive notifications triggered when trend
    analysis suggests a potential threshold breach.
  \item \textbf{Infrastructure Notifications:} Continuous monitoring of hardware health,
    evaluating Thread node reachability and gateway connectivity.
\end{enumerate}


% ============================================================
\section{Hazard Analysis \& Reliability}

Given the context of environmental monitoring, the system implements mitigation
strategies for operational and distributed system hazards:

\begin{table}[htbp]
\centering
\setlength{\tabcolsep}{4pt}
\renewcommand{\arraystretch}{1.0}
\caption{Hazard Analysis \& Reliability Mitigation}
\label{tab:hazards}
\small
\begin{tabularx}{\textwidth}{@{} >{\raggedright\arraybackslash}p{3.8cm} >{\raggedright\arraybackslash}X @{}}
\toprule
\textbf{Hazard / Failure Mode} & \textbf{Mitigation / Reliability Strategy} \\
\midrule
ML Calibration Phase
  & Deterministic safety limits remain active independently during model warm-up. \\
\midrule
Model Drift (False Positives)
  & Isolation Forest utilizes a rolling 24-hour window, automatically absorbing slow environmental drift. \\
\midrule
Network Loss (No ISP)
  & Fog-Native architecture allows local LAN users to receive alerts even without external internet. \\
\midrule
CoAP/UDP Packet Loss
  & CoAP \texttt{CONFIRMABLE} messages enforce acknowledgments, triggering retransmissions to reduce data gaps. \\
\midrule
Mesh Node Failure
  & Thread multipath routing bypasses dead nodes. Isolated nodes utilize onboard flash for temporary Store-and-Forward buffering. \\
\midrule
Sensor Drift / Failure
  & Plausibility validation, stuck-value detection, and cross-sensor consistency checks. \\
\midrule
Gateway / Service Crash
  & Redundant Docker restart policies, heartbeat monitoring, and decoupled deterministic safety thresholds. \\
\bottomrule
\end{tabularx}
\end{table}

% ============================================================
\section{Evaluation Methodology}

To validate the proposed system, the following evaluation metrics will be measured
during the testing phase:

\begin{itemize}
  \item \textbf{Network Performance:} Thread mesh reliability, telemetry latency, and
    packet delivery ratio under simulated indoor interference.
  \item \textbf{ML Effectiveness:} Anomaly detection accuracy, false positive rate, and
    the convergence speed of the Isolation Forest during the calibration phase.
  \item \textbf{System Scalability \& Resource Utilization:} CPU/Memory footprint of
    the Gateway-Tier Fog microservices and response times of the Spring Boot API under
    increased simulated node loads.
\end{itemize}

% ============================================================
\section{Limitations \& Expected Outcomes}

\textbf{Limitations:}
\begin{itemize}
  \item \textbf{Deployment Scale:} Current physical validation is constrained to a
    localized indoor environment utilizing a single physical sensor node. The proposed
    architecture is designed to support scalable multi-node OpenThread mesh deployments.
  \item \textbf{Training Duration:} The rolling baseline relies on a 24-hour window,
    meaning long-term seasonal environmental changes will not be captured in the
    immediate dataset.
  \item \textbf{Hardware Constraints:} The lightweight nature of the ML models restricts
    the implementation of highly complex, multi-layered predictive neural networks.
\end{itemize}

\textbf{Expected Outcomes:}
The project aims to deliver a fully functional, end-to-end prototype that demonstrates
the viability of combining low-power Thread mesh networking with localized, edge-native
anomaly detection. The resulting architecture will provide a reproducible framework for
privacy-preserving, intelligent indoor monitoring.

% ============================================================
\section{Project Management Summary}

\subsection{Project Team}

\begin{table}[htbp]
\centering
\setlength{\tabcolsep}{4pt}
\renewcommand{\arraystretch}{1.0}
\caption{Project Team Structure}
\label{tab:team}
\small
\begin{tabularx}{\textwidth}{@{} >{\raggedright\arraybackslash}p{3.5cm} >{\raggedright\arraybackslash}p{4cm} >{\raggedright\arraybackslash}X @{}}
\toprule
\textbf{Team Members} & \textbf{Role} & \textbf{Responsibilities} \\
\midrule
M.~F.~Mufti / M.~Anwar
  & Embedded Systems Engineer
  & OpenThread firmware, sensor integration, CoAP communication \\
\midrule
M.~I.~Hussain
  & Machine Learning Engineer
  & Anomaly detection, trend analysis, XAI integration \\
\midrule
H.~T.~Khan / M.~Saleem
  & Backend \& Database Engineer
  & Ingestion API, database schema design, persistence \\
\midrule
H.~T.~Khan / S.~A.~Minhas
  & Frontend Developer
  & Dashboard UI/UX, visualization, frontend integration \\
\midrule
M.~Anwar / M.~Saleem
  & Gateway \& Infrastructure Engineer
  & Serial Gateway development, Docker orchestration \\
\midrule
M.~Anwar / S.~A.~Minhas
  & Scrum Master \& Documentation Lead
  & Sprint planning, technical documentation, presentations \\
\bottomrule
\end{tabularx}
\end{table}

\newpage
\subsection{Sprint Schedule}

The project utilizes an Agile framework with one-week sprints over a 10-week lifecycle, executing cross-functional development tracks concurrently:

\begin{table}[htbp]
\centering
\setlength{\tabcolsep}{4pt}
\renewcommand{\arraystretch}{1.0}
\caption{Agile Sprint Schedule}
\label{tab:sprint}
\small
\begin{tabularx}{\textwidth}{@{} >{\bfseries}p{2cm} >{\raggedright\arraybackslash}X @{}}
\toprule
Sprint & Concurrent Cross-Functional Tracks \\
\midrule
Sprint 1--2
  & OpenThread sensor provisioning (Embedded) \& database schema design (Backend) \\
\midrule
Sprint 3--4
  & Serial Gateway bridging (Infrastructure) \& React UI/UX development (Frontend) \\
\midrule
Sprint 5--6
  & CoAP-to-HTTP telemetry ingestion (Backend) \& Isolation Forest / Trend Analysis (ML) \\
\midrule
Sprint 7--8
  & Explainable AI alert integration (ML) \& real-time Dashboard data binding (Frontend) \\
\midrule
Sprint 9--10
  & Docker orchestration, system calibration, and academic documentation (Cross-Functional) \\
\bottomrule
\end{tabularx}
\end{table}

\subsection{Workload Estimation}

A hybrid approach combining Bottom-Up Estimation and Complexity-Based Estimation was adopted, yielding a total workload of approximately 600 cumulative team hours over the 10-week lifecycle. Subsystem tasks were individually estimated and refined based on technical complexity, with additional effort allocated to high-risk components including Thread/CoAP network validation, ML baseline calibration, and Docker multi-service orchestration.

\begin{table}[htbp]
\centering
\setlength{\tabcolsep}{4pt}
\renewcommand{\arraystretch}{1.0}
\caption{Estimation Methods}
\label{tab:estimation}
\small
\begin{tabularx}{\textwidth}{@{} >{\raggedright\arraybackslash}p{4cm} >{\raggedright\arraybackslash}X @{}}
\toprule
\textbf{Estimation Method} & \textbf{Purpose} \\
\midrule
Bottom-Up Estimation
  & Subsystem-level workload estimation and task allocation \\
\midrule
Complexity-Based Estimation
  & Integration, technical risk, and complexity evaluation \\
\bottomrule
\end{tabularx}
\end{table}

% ============================================================
\vspace{-0.5em}
\begin{thebibliography}{8}
\setlength{\itemsep}{-2pt}
\setlength{\parskip}{0pt}

\bibitem{ortega2025}
E.~Ortega, F.~Su, R.~Chattopadhyay, and K.~Chakrabarty,
``Discretized-Isolation Forest: Memory- and Compute-Efficient Unsupervised Anomaly
Detection for Resource-Constrained Internet of Things Edge Devices,''
\textit{IEEE Internet of Things Journal}, vol.~12, no.~2, pp.~1699--1717, Jan.\ 2025.
\newblock \href{https://doi.org/10.1109/JIOT.2024.3468950}{doi:10.1109/JIOT.2024.3468950}

\bibitem{wu2022}
Y.~Wu, H.-N.~Dai, and H.~Tang,
``Graph Neural Networks for Anomaly Detection in Industrial Internet of Things,''
\textit{IEEE Internet of Things Journal}, vol.~9, no.~12, pp.~9214--9231, Jun.\ 2022.
\newblock \href{https://doi.org/10.1109/JIOT.2021.3094295}{doi:10.1109/JIOT.2021.3094295}

\bibitem{chaturvedi2026}
P.~Chaturvedi, S.~Ahmad, and A.~Mewada,
``Fog Computing for Intelligent and Sustainable IoT: Architectures, Optimization,
Datasets, and Open Challenges,''
in \textit{Proc.\ 2026 IEEE Madhya Pradesh Section Conference (MPCON)}, Gwalior, India, 2026.
\newblock \href{https://doi.org/10.1109/MPCON69668.2026.11508383}{doi:10.1109/MPCON69668.2026.11508383}

\bibitem{chhabra2025}
G.~S.~Chhabra, G.~N.~V.~Rajareddy, A.~Mahapatra, S.~S.~Mangalampalli, K.~S.~Sahoo, D.~Sethi, and K.~Mishra,
``Deep Learning-Centric Task Offloading in IoT-Fog-Cloud Continuum: A State-of-the-Art
Review, Open Research Issues, and Future Directions,''
\textit{IEEE Access}, vol.~13, pp.~144241--144270, Aug.\ 2025.
\newblock \href{https://doi.org/10.1109/ACCESS.2025.3599190}{doi:10.1109/ACCESS.2025.3599190}

\bibitem{guha2025}
S.~Guha and A.~Datta,
``GDRNet: A Deep Learning Framework for Dynamic Sliding Window Determination and
Improved Prediction Accuracy in Online IIoT Systems,''
\textit{IEEE Access}, vol.~13, 2025.
\newblock \href{https://doi.org/10.1109/ACCESS.2025.3634671}{doi:10.1109/ACCESS.2025.3634671}

\bibitem{who2021}
World Health Organization,
\textit{WHO Global Air Quality Guidelines: Particulate Matter (PM\textsubscript{2.5}
and PM\textsubscript{10}), Ozone, Nitrogen Dioxide, Sulfur Dioxide and Carbon Monoxide}.
Geneva: WHO Press, 2021.
\newblock \href{https://www.who.int/publications/i/item/9789240034228}{Available online}

\bibitem{khattak2023}
S.~B.~A.~Khattak, M.~M.~Nasralla, H.~Farman, and N.~Choudhury,
``Performance Evaluation of an IEEE 802.15.4-Based Thread Network for Efficient Internet
of Things Communications in Smart Cities,''
\textit{Applied Sciences (MDPI)}, vol.~13, no.~13, p.~7745, 2023.
\newblock \href{https://doi.org/10.3390/app13137745}{doi:10.3390/app13137745}

\bibitem{coap2014}
Z.~Shelby, K.~Hartke, and C.~Bormann,
``The Constrained Application Protocol (CoAP),''
\textit{IETF RFC~7252}, Internet Engineering Task Force, Jun.\ 2014.
\newblock \href{https://www.rfc-editor.org/rfc/rfc7252}{Available online}

\end{thebibliography}

\end{document}
